% 将军饮马模型主图：直线 l, 同侧两点 A, B；A 的对称点 A'；最优 P 为 A'B 与 l 的交点
\documentclass[tikz,border=4pt]{standalone}
\usepackage{ctex}
\usepackage{amsmath}
\usetikzlibrary{calc, arrows.meta}
\begin{document}
\begin{tikzpicture}[scale=1.0, thick]
  % 直线 l：水平
  \draw[thick] (-0.5, 0) -- (7.5, 0);
  \node[right] at (7.5, 0) {$l$};

  % A, B 在 l 上方
  \coordinate (A) at (1.5, 2);
  \coordinate (B) at (5.5, 3);
  % A 关于 l 的对称点 A'
  \coordinate (Ap) at (1.5, -2);

  % A'B 与 l 的交点 P：参数化 (1-t)(1.5,-2) + t(5.5,3) = (1.5+4t, -2+5t)；y=0 ⇒ t=2/5 ⇒ x=1.5+1.6=3.1
  \coordinate (P) at (3.1, 0);

  \filldraw (A) circle (2pt) node[above] {$A$};
  \filldraw (B) circle (2pt) node[above] {$B$};
  \filldraw (Ap) circle (2pt) node[below] {$A'$};
  \filldraw (P) circle (2pt) node[below right] {$P$};

  % 折线 PA + PB（红实线）
  \draw[red, thick] (A) -- (P) -- (B);
  % PA'（虚线，等于 PA）
  \draw[blue, dashed, thick] (P) -- (Ap);
  % A 与 A' 的连线（中垂）
  \draw[gray, dotted] (A) -- (Ap);
  % 中垂直角符号
  \draw (1.5, 0)--(1.7, 0)--(1.7,-0.2)--(1.5,-0.2);

  \node[red, font=\footnotesize] at (2.2, 1.2) {$PA$};
  \node[red, font=\footnotesize] at (4.5, 1.8) {$PB$};
  \node[blue, font=\footnotesize] at (2.0, -1.2) {$PA'$};

  \node[font=\footnotesize, gray] at (4.6, -0.3) {$A'B$ 与 $l$ 交点即最优 $P$};
\end{tikzpicture}
\end{document}
